% §6 Experimental Setup — target 0.75 page

\section{Experimental Setup}
\label{sec:setup}

\paragraph{Infrastructure}
All scenarios run on a single Proxmox VE host. Networks are Linux bridges
(\texttt{vmbr1}, \texttt{192.168.100.0/24}) with an OpenWrt router as gateway.
Per-scenario deployment and teardown are automated via Ansible and complete
in under 10~minutes per scenario.

\paragraph{LLM model — fairness lock}
To isolate \emph{architectural} contributions from model-specific advantages, all
LLM-driven systems compared in this work --- our pipeline (and its ablations) and
all LLM baselines --- run on the \emph{same} general-purpose model:
\textbf{MiniMax-M2.7} \cite{minimax27}. We deliberately avoid a multi-model
matrix: a fair comparison of architectures requires the model variable to be
held constant. CAI's specialised \texttt{alias1} model (security-fine-tuned,
paywalled) is intentionally not used --- this study addresses the architecture
axis, not the model fine-tuning axis. The choice of MiniMax-M2.7 also keeps the
combined API cost at zero under our research subscription, enabling the full
$k{=}3$-run variance estimation across all systems and scenarios without budget
trade-offs.

\paragraph{Baselines}
We compare our pipeline against three categories of systems, each occupying a
distinct architectural slot:

\begin{itemize}
  \item \textbf{Single-agent LLM} --- \textsc{PentestGPT}~\cite{pentestgpt24}
  in autonomous mode, the canonical USENIX'24 reference.
  \item \textbf{Multi-agent LLM (no infrastructure graph)} ---
  \textsc{CAI}~\cite{cai25}. We test two variants: \emph{A1} per-IP
  (one independent CAI session per ground-truth IP, no shared state) and
  \emph{B} CIDR-scope (a single CAI session given the full \texttt{/24}, allowed
  to discover hosts and pivot via \texttt{LinuxCmd} primitives). Both variants
  receive the same total turn budget as our pipeline.
  \item \textbf{Multi-agent LLM with task-graph} ---
  \textsc{VulnBot}~\cite{vulnbot25}, the most structurally similar competitor.
  VulnBot's Penetration Task Graph (PTG) is a directed acyclic graph of
  \emph{tasks} (action dependencies), to be contrasted with our infrastructure
  graph of \emph{devices} (firewall-enforced reachability). Both are
  ``graph-based'', but encode orthogonal abstractions.
  \item \textbf{Non-LLM scanner} --- Nmap with all NSE scripts
  (\texttt{--script=default,vuln,auth}) against the same target CIDR. Scanner
  outputs are converted to the evaluator's JSON schema.
\end{itemize}

\paragraph{Internal ablations}
Two ablations isolate the contribution of our two key design choices:
(i)~\harnessname--w/o-Phase~1 disables graph analysis (\texttt{--phases 2 3 4 5
6}), testing the value of infrastructure-graph guidance;
(ii)~\harnessname--w/o-Phase~5 disables the intrusion phase (\texttt{--phases 1
2 3 4 6}), testing the value of structured multi-hop lateral movement.

\paragraph{Variance}
Each (system, scenario) cell is run $k{=}3$ times. We report mean, standard
deviation, and 95\% bootstrap confidence intervals. Pairwise differences
(our pipeline vs each baseline) are tested with Mann--Whitney $U$ on F1 and
weighted Score per scenario. This protocol exceeds the single-run reporting of
PentestGPT~\cite{pentestgpt24}, AutoPentester~\cite{autopentester25},
AUTOPENBENCH~\cite{autopenbench24}, and CAI~\cite{cai25}, and matches the
multi-run averaging of VulnBot~\cite{vulnbot25} (5-experiment) and
Pentest-R1~\cite{pentestr125} (pass@3).

\paragraph{Metrics}
We report five families of metrics, chosen to (i)~align with prior work's
expected vocabulary and (ii)~surface architectural differences invisible to
host-level benchmarks:
\begin{itemize}
  \item \emph{Task completion} --- Recall against the ground-truth vulnerability
  set, the network-scale analogue of the sub-task / task completion rate
  reported by PentestGPT, VulnBot, AutoPentester, Pentest-R1, and xOffense.
  \item \emph{Quality} --- Precision, F1, and weighted Score
  (CVSS-severity-weighted Recall, \S\ref{sec:iotbench:eval}). Score plays the
  role of AutoPentester's ``vulnerability coverage'' but with severity weights.
  \item \emph{Multi-Hop Reach} --- MHR$_k$ for $k\in\{1,2,3\}$, the fraction of
  ground-truth vulnerabilities at \texttt{hop\_depth}~$\geq k$ that are reached.
  No prior LLM-pentest benchmark measures network depth explicitly.
  \item \emph{Evidence depth} --- $L_1$ (detection), $L_2$ (exploitation), $L_3$
  (exfiltration) coverage; absent from binary flag-capture benchmarks.
  \item \emph{Cost envelope} --- prompt+completion tokens and monetary cost
  (USD) from provider billing APIs, plus wall-clock time per scenario,
  matching CAI~\cite{cai25}, AutoPentester, and Pentest-R1.
\end{itemize}

\paragraph{Compute and budget}
All runs consume a combined API budget under \$NNN, bounded by per-phase turn
caps and the concurrency limit in Phase~4.
\todo{Insert exact budget after all runs complete (freeze 2 May).}
